\section{Problem Formulation}
\label{sec:problem}
We retain the prelaunch allocation problem introduced in
WORKS~\cite{camp2026works}. A workflow comprises named processes, each of
which instantiates tasks for different samples or input partitions.
Before task $t$ starts, the allocator issues an initial memory request
$\hat y_t$. Its static feature vector $\mathbf{s}_t$ contains workflow and
process metadata, software and system context, resource settings, and
staged input size. Execution subsequently reveals peak resident set size
(RSS), $y_t$, and duration $d_t$. If the task is audited, it also supplies a
consumption measurement $c_t$. Neither $y_t$ nor $c_t$ is available for that
task's initial allocation.

Let $\mathcal{P}(t)$ contain tasks whose outcomes have been recorded before
the allocation decision for $t$. Peak-memory history $\mathbf{h}_{m,t}$
summarizes compatible observations from this set; consumption history
$\mathbf{h}_{c,t}$ uses its audited subset. The prelaunch decisions are
\begin{equation}
\hat y_t=f(\mathbf{s}_t,\mathbf{h}_{m,t},\mathbf{h}_{c,t}),
\qquad q_t\in\{0,1\},
\label{eq:allocation_problem}
\end{equation}
where $f$ is the memory-allocation function and $q_t=1$ selects auditing.
A budget $B\in(0,1]$ limits the audited
fraction over the specified evaluation horizon. Auditing task $t$ cannot
improve the request already issued to it; its measurements can inform
only later decisions. During live execution, these results become
available only after the task finishes and they are added to CAMP's
history.

The allocation objective is to limit underallocation while avoiding
unnecessary reservations. For $N$ evaluated tasks, coverage is the
fraction whose initial request meets or exceeds measured peak memory:
\begin{equation}
\mathrm{Cov}=\frac{1}{N}\sum_{t=1}^{N}\mathbf{1}[\hat y_t\geq y_t],
\label{eq:coverage}
\end{equation}
where $\mathbf{1}[\cdot]$ is the indicator function. As in WORKS, retry
increases do not enter $\hat y_t$. Live studies therefore
report task completion, failures, and retries separately.

We evaluate coverage alongside total requested memory, $\sum_t\hat y_t$
(GiB), and unused memory-time, $W$ (GiB-hours), defined as follows:
\begin{equation}
W=\sum_{t\in\mathcal{D}}\max(\hat y_t-y_t,0)\,d_t,
\label{eq:unused_memory_time}
\end{equation}
where $\mathcal{D}$ contains tasks with valid runtime measurements, memory
is expressed in GiB, and runtime in hours. For each task, $W$ multiplies
memory reserved above its measured peak by its runtime, then sums these
values across tasks. It uses peak memory rather than tracking memory
usage throughout execution. Total requested memory sums requests across
tasks that may run at different times; it does not measure memory reserved
simultaneously. The live experiments report $W$; a separate audit-cost
analysis estimates the extra reserved memory-time caused by auditing,
also in GiB-hours. Section~\ref{sec:experiments} defines this auditing cost
and specifies the runtime and cost components included in each analysis.

To evaluate whether the additional warm run improves allocation on unseen
inputs, we compare two paths from the same cold-run checkpoint:
\[
\begin{aligned}
\text{Warm-Transfer (WT):}\quad &\text{Cold} \rightarrow \text{Warm} \rightarrow \text{New inputs},\\
\text{Direct-Transfer (DT):}\quad &\text{Cold} \rightarrow \text{New inputs}.
\end{aligned}
\]
Warm repeats the original inputs; both paths then process the same unseen
inputs. Cold is shared, not rerun for each path. Both branches update their
models, history, and calibration online, excluding a task's own outcomes
from its prelaunch request. We assess the additional warm run using
coverage, total requested memory, and unused memory-time.
